% GENERATED FILE. Edit canonical lessons, then run npm run generate:books.
% canonical-source-hash: fnv1a64:3591a8a010f882ff
% canonical-lessons: BN-C509-galicha, BN-C509-katha, BN-C509-paposh, BN-C509-godi, BN-C509-bichhano

\chapter{Rugs and Quilts}
\label{ch:bn-rugs-and-quilts}
\hlchaptermodality{509}

\begin{chapteropening}
\textbf{By the end of this chapter:} \emph{I can talk about rugs, quilts and cushions and spreading things out.}

Complete the last lesson of chapter 509: I can talk about rugs, quilts and cushions and spreading things out.
\end{chapteropening}

\section[\texorpdfstring{\bn{গালিচা} (gālichā)}{গালিচা (gālichā)}]{\bn{গালিচা} (gālichā) — a rug; a carpet}

\label{lesson:BN-C509-galicha}



\begin{quote}
Before the new one: say the Bengali for a door frame, then the Bengali for to hang up.
\end{quote}

\subsection*{You'll want to know: \bn{গালিচা}}
\textbf{\bn{গালিচা}} — \emph{gālichā} — \textquotedblleft{}a rug; a carpet\textquotedblright{}.

\begin{quote}
\bn{বসার} \bn{ঘরে} \bn{একটা} \bn{লাল} \bn{গালিচা} \bn{আছে।}
\end{quote}

\begin{quote}
\emph{Bôsār ghôre ekṭā lāl gālichā āchhe.}
\end{quote}

\begin{quote}
There's a red rug in the sitting room.
\end{quote}

\begin{cousinweb}
\textbf{\bn{গালিচা}} (\emph{gālichā}) is usually traced to Persian \emph{ghālīcha}, \textquotedblleft{}a small carpet\textquotedblright{}, itself thought to come from Turkish.
\end{cousinweb}

\begin{grammarlens}[title={What you've built}]
One more word for describing what you lived and felt.
\end{grammarlens}

\subsection*{Your turn}
\begin{itemize}
\raggedright
  \item \emph{Say it:} \emph{gālichā}
  \item \emph{Say it:} \emph{gālichā}, once more
  \item \emph{Say it:} \emph{Bôsār ghôre ekṭā lāl gālichā āchhe.} — There's a red rug in the sitting room.
  \item \emph{Recall:} say the Bengali for a door frame, then the Bengali for to hang up, then say \emph{gālichā} again
\end{itemize}

\begin{tcolorbox}[breakable,colback=teal!4,colframe=teal!35!black,title={Before you move on}]
What does \bn{গালিচা} mean? (\textquotedblleft{}A rug; a carpet\textquotedblright{}.) Say it once more.
\end{tcolorbox}
\section[\texorpdfstring{\bn{কাঁথা} (kā̃thā)}{কাঁথা (kā̃thā)}]{\bn{কাঁথা} (kā̃thā) — a kantha (an embroidered quilt)}

\label{lesson:BN-C509-katha}



\begin{quote}
Before the new one: say the Bengali for to hang up, then the Bengali for a rug.
\end{quote}

\subsection*{You'll want to know: \bn{কাঁথা}}
\textbf{\bn{কাঁথা}} — \emph{kā̃thā} — \textquotedblleft{}a kantha (an embroidered quilt)\textquotedblright{}.

\begin{quote}
\bn{শীতে} \bn{দিদিমা} \bn{আমাকে} \bn{একটা} \bn{কাঁথা} \bn{দিলেন।}
\end{quote}

\begin{quote}
\emph{Shīte didimā āmāke ekṭā kā̃thā dilen.}
\end{quote}

\begin{quote}
In winter my grandmother gave me a kantha quilt.
\end{quote}

\begin{cousinweb}
\textbf{\bn{কাঁথা}} (\emph{kā̃thā}) goes back to Sanskrit \emph{kanthā}, \textquotedblleft{}a patched garment, a rag\textquotedblright{} (Monier-Williams): old cloth stitched into new use.
\end{cousinweb}

\begin{grammarlens}[title={What you've built}]
One more word for describing what you lived and felt.
\end{grammarlens}

\subsection*{Your turn}
\begin{itemize}
\raggedright
  \item \emph{Say it:} \emph{kā̃thā}
  \item \emph{Say it:} \emph{kā̃thā}, once more
  \item \emph{Say it:} \emph{Shīte didimā āmāke ekṭā kā̃thā dilen.} — In winter my grandmother gave me a kantha quilt.
  \item \emph{Recall:} say the Bengali for to hang up, then the Bengali for a rug, then say \emph{kā̃thā} again
\end{itemize}

\begin{tcolorbox}[breakable,colback=teal!4,colframe=teal!35!black,title={Before you move on}]
What does \bn{কাঁথা} mean? (\textquotedblleft{}A kantha (an embroidered quilt)\textquotedblright{}.) Say it once more.
\end{tcolorbox}
\section[\texorpdfstring{\bn{পাপোশ} (pāposh)}{পাপোশ (pāposh)}]{\bn{পাপোশ} (pāposh) — a doormat}

\label{lesson:BN-C509-paposh}



\begin{quote}
Before the new one: say the Bengali for a rug, then the Bengali for an embroidered quilt.
\end{quote}

\subsection*{You'll want to know: \bn{পাপোশ}}
\textbf{\bn{পাপোশ}} — \emph{pāposh} — \textquotedblleft{}a doormat\textquotedblright{}.

\begin{quote}
\bn{দরজার} \bn{সামনে} \bn{একটা} \bn{পাপোশ} \bn{আছে।}
\end{quote}

\begin{quote}
\emph{Dôrjār sāmne ekṭā pāposh āchhe.}
\end{quote}

\begin{quote}
There's a doormat in front of the door.
\end{quote}

\begin{cousinweb}
\textbf{\bn{পাপোশ}} (\emph{pāposh}) is usually traced to Persian \emph{pāposh}, \textquotedblleft{}a foot covering\textquotedblright{}: \emph{pā}, \textquotedblleft{}foot\textquotedblright{}, and \emph{posh}, \textquotedblleft{}covering\textquotedblright{}.
\end{cousinweb}

\begin{grammarlens}[title={What you've built}]
One more word for describing what you lived and felt.
\end{grammarlens}

\subsection*{Your turn}
\begin{itemize}
\raggedright
  \item \emph{Say it:} \emph{pāposh}
  \item \emph{Say it:} \emph{pāposh}, once more
  \item \emph{Say it:} \emph{Dôrjār sāmne ekṭā pāposh āchhe.} — There's a doormat in front of the door.
  \item \emph{Recall:} say the Bengali for a rug, then the Bengali for an embroidered quilt, then say \emph{pāposh} again
\end{itemize}

\begin{tcolorbox}[breakable,colback=teal!4,colframe=teal!35!black,title={Before you move on}]
What does \bn{পাপোশ} mean? (\textquotedblleft{}A doormat\textquotedblright{}.) Say it once more.
\end{tcolorbox}
\section[\texorpdfstring{\bn{গদি} (gôdi)}{গদি (gôdi)}]{\bn{গদি} (gôdi) — a cushion; a padded seat}

\label{lesson:BN-C509-godi}



\begin{quote}
Before the new one: say the Bengali for an embroidered quilt, then the Bengali for a doormat.
\end{quote}

\subsection*{You'll want to know: \bn{গদি}}
\textbf{\bn{গদি}} — \emph{gôdi} — \textquotedblleft{}a cushion; a padded seat\textquotedblright{}.

\begin{quote}
\bn{এই} \bn{চেয়ারের} \bn{গদিটা} \bn{নরম।}
\end{quote}

\begin{quote}
\emph{Ei cheyārer gôdiṭā nôrôm.}
\end{quote}

\begin{quote}
This chair's cushion is soft.
\end{quote}

\begin{cousinweb}
\textbf{\bn{গদি}} (\emph{gôdi}) is usually traced to Hindi \emph{gaddī}, \textquotedblleft{}a cushion, a padded seat\textquotedblright{}; its earlier origin is not settled.
\end{cousinweb}

\begin{grammarlens}[title={What you've built}]
One more word for describing what you lived and felt.
\end{grammarlens}

\subsection*{Your turn}
\begin{itemize}
\raggedright
  \item \emph{Say it:} \emph{gôdi}
  \item \emph{Say it:} \emph{gôdi}, once more
  \item \emph{Say it:} \emph{Ei cheyārer gôdiṭā nôrôm.} — This chair's cushion is soft.
  \item \emph{Recall:} say the Bengali for an embroidered quilt, then the Bengali for a doormat, then say \emph{gôdi} again
\end{itemize}

\begin{tcolorbox}[breakable,colback=teal!4,colframe=teal!35!black,title={Before you move on}]
What does \bn{গদি} mean? (\textquotedblleft{}A cushion; a padded seat\textquotedblright{}.) Say it once more.
\end{tcolorbox}
\section[\texorpdfstring{\bn{বিছানো} (bichhāno)}{বিছানো (bichhāno)}]{\bn{বিছানো} (bichhāno) — to spread out (a mat or a sheet)}

\label{lesson:BN-C509-bichhano}



\begin{quote}
Before the new one: say the Bengali for a doormat, then the Bengali for a cushion.
\end{quote}

\subsection*{You'll want to know: \bn{বিছানো}}
\textbf{\bn{বিছানো}} — \emph{bichhāno} — \textquotedblleft{}to spread out (a mat or a sheet)\textquotedblright{}.

\begin{quote}
\bn{আমরা} \bn{মেঝেতে} \bn{মাদুর} \bn{বিছিয়ে} \bn{খেলাম।}
\end{quote}

\begin{quote}
\emph{Āmrā mejhete mādur bichhiye khelām.}
\end{quote}

\begin{quote}
We spread a mat on the floor and ate.
\end{quote}

Past (he/she …): \bn{তিনি} \bn{বিছানায়} \bn{চাদর} \bn{বিছালেন}, \emph{tini bichhānāy chādôr bichhālen}.

\begin{cousinweb}
\textbf{\bn{বিছানো}} (\emph{bichhāno}) is a tadbhava from a reconstructed \emph{vicchādayati}; CDIAL lists Bengali \emph{bichāna}, \textquotedblleft{}to spread, to strew\textquotedblright{} (CDIAL 11692). \bn{বিছানা}, a bed, comes from the same family.
\end{cousinweb}

\begin{grammarlens}[title={What you've built}]
One more word for describing what you lived and felt.
\end{grammarlens}

\subsection*{Your turn}
\begin{itemize}
\raggedright
  \item \emph{Say it:} \emph{bichhāno}
  \item \emph{Say it:} \emph{bichhāno}, once more
  \item \emph{Say it:} \emph{Āmrā mejhete mādur bichhiye khelām.} — We spread a mat on the floor and ate.
  \item \emph{Recall:} say the Bengali for a doormat, then the Bengali for a cushion, then say \emph{bichhāno} again
\end{itemize}

\begin{tcolorbox}[breakable,colback=teal!4,colframe=teal!35!black,title={Before you move on}]
What does \bn{বিছানো} mean? (\textquotedblleft{}To spread out (a mat or a sheet)\textquotedblright{}.) Say it once more.
\end{tcolorbox}
